\chapter{Generating Traces}\label{ch:traces}
In this section we describe the process of generating game traces then transforming them into to IGGP tasks. The first step is to generate the intelligent and random gameplay. To do this we used the Sancho system \cite{Sancho/Github}.

\section{Sancho}\label{sec:sancho}
Generating intelligent gameplay is a nebulous task. Defining intelligence is no small ask. Using optimal gameplay traces or human generated gameplay would be ideal, however this is infeasible in the scope of this project. Provably optimal traces were experimented with. An optimal player for \textit{eight puzzle} using A* with an admissible heuristic was created with the intention of writing optimal players for more games such as \textit{tic tac toe} and \textit{knights tour}. However when testing the Sancho system it was found that game traces produced were often of near optimal quality. In fact for \textit{eight puzzle} it generates provably optimal traces as described later in the chapter. It was decided that it would be easier and we would generate a possibly more a consistent level of play across all games using the GGP player Sancho.

The winner of the 2014 GGP competition ``Sancho''\footnote{http://ggp.stanford.edu/iggpc/winners.php accessed on 12/03/2020} was selected since it is the most recent winner with published code \cite{Sancho/Github}. Some small modifications to the information logged by the game server are made but otherwise the code is unchanged. Since the aim of the GGP competition is to find the program that performs best at the set of games used in the IGGP problem we conclude that this system was the best general game player to use. We can assume Sancho can play above amateur human level across the majority of games in the the dataset. At the annual GGP competition there is also a match held at the end where the human creators of the AI themselves play against the winner. The humans almost never win \cite{Genesereth/GGPOverview}.

Sancho uses a range of techniques to generate gameplay. The core of algorithm used by Sancho is the Monte Carlo tree search (MCTS).

\subsection{Monte Carlo Tree Search}
Given a game state the basic MCTS will return the most promising next move. The algorithm achieves this by simulating random playouts of the game many times. The technique was developed for Computer Go but has since been applied to play a wide range of games effectively including board games and video games \cite{Silver/MCTS,Chaslot/MCTS}. The use of random simulation to evaluate game states is a powerful tool. No extra information about the game such as heuristics for evaluating states are needed at all, the rules of the game suffice. In the case of the GGP problem this is ideal.

In our case all games being played are sequential, finite and discrete so we only need to consider MCTS for this case.

\begin{figure}[ht]
	\centering
	\includegraphics[width=0.9\linewidth]{Tic-tac-toe-game-tree.png}
	\caption{A section of a game tree of \textit{tic tac toe} showing some of the possible moves. \textit{Stannered/CC-BY-SA-3.0}}
	\label{fig:game-tree}
\end{figure}

The MCTS is a tree search algorithm, the tree being searched is the game tree. Each node on the game tree represents a state (figure \ref{fig:game-tree}). The search is a sequence of traversals of the game tree. A \textit{traversal} is a path that starts at the root node and continues down until it reaches a node that has at least one unvisited child (not necessarily a terminal state). One of these unvisited children is then chosen to be the start state for a simulation of the rest of the game. The simulation selects moves randomly, playing the game out to a terminal state. The result of the simulation is propagated back from the node it started at all the way to the root node, updating statistics attached to each node. These statistics are used to choose future paths to traverse so that more promising moves are investigated with higher probability.

\subsubsection{Use of MCTS in Sancho}

Sancho makes a few modifications to MCTS\footnote{https://sanchoggp.blogspot.com/2014/05/what-is-sancho.html accessed on 15/03/2020}. The main one being the adding of heuristics. It also makes optimisations to increase efficiency.

In GGP matches a period of time before the match is given in which to do pre match calculations. Sancho uses this period to derive basic heuristics. A heuristic should take a game state or move and return a value based on how promising that state or move is in relation to the goal state. 

The identification of possible heuristics takes place in two stages. The first of these is static analysis of the game rules. This static analysis identifies possible heuristics that can be applied to the current game. These include things like piece capture: if certain rules indicate capture of the players piece these can be selected against. Another is numeric quantity detection: a number in the state acts as a heuristic (like number of coins a player has). The second stage is simulation of the game. Many (possibly tens of thousands) of full simulations of the game are run. After the simulations are complete a correlation coefficient is calculated between each candidate heuristic's observed values and the eventual score achieved in the game. Those heuristics that show correlation above a fixed threshold are then enabled, and will be used during play to guide the tree exploration in MCTS.

\subsection{Other aspects of Sancho}

Monte Carlo simulations are used for most of the games available however Sancho identifies some games that can be solved more efficiently in other ways. Any single player game in the GGP dataset can be analysed using puzzle solving techniques.
\subsubsection{Puzzle Solving}
A single player game in the GGP dataset is necessarily a deterministic game of complete information due to the constraints of the Game Description Language. This gives the useful constraint that any solution found in a playout of the game will always remain valid since the game is completely deterministic. Sancho identifies these single player games and attempts to derive heuristics as described above. Where an obvious goal state exists a distance metric such as the hamming distance between two states can be applied. A* search is then used to find optimal solutions. For some games, such as \textit{eight puzzle} the hamming distance is an admissible heuristic so Sancho finds a provably optimal solution.

\subsubsection{Static analysis}

Sancho also does static analysis of the game rules. It determines game predicates that if true imply that either the goal will never be reached or the goal will always be reached. An example of this is a game such as \textit{untwisty corridor} (the game is effectively a maze where you immediately lose if you step off a ``safe'' path). If the safe path is stepped off then the goal can never be reached so any state not on the safe path is avoided.

\subsection{Sancho and IGGP}
Sancho generally provides good quality intelligent play when run on the games in the IGGP dataset however in some cases it struggles.

The dataset contains several games based on game theory. Games such as the \textit{prisoner dilemma} where player can choose to either \textit{cooperate} or \textit{defect}. If both players cooperate then they both get 3 points. If one defects and the other tried to cooperate the cooperator gets 0 points and the defector gets 5. If they both defect both get 0. There are several rounds of this in one game playout. A human playing this might choose to try and cooperate, hoping that that the other player will too however Sancho has been observed to always defect. This result is predictable since this is the only strong Nash equilibrium for this game however it is hard to justify this as ``human quality play''.


\section{Generating and transforming the traces}\label{sec:gen}

To generate the intelligent and random traces for the experiments an automatic system was constructed. The system runs Sancho or a random game player on the GDL game descriptions and converts the traces into IGGP tasks.

The general game playing community have developed a codebase that provides the basic functionality for hosting a match between GGP agents\footnote{https://github.com/ggp-org/ggp-base accessed on 13/03/2020}. This codebase along with the GGP agent Sancho is used to generate the datasets. The codebase includes a very basic random GGP player that will, given a GDL game description, play a random legal move every turn. To generate the random traces for a game of $x$ players $x$ Sancho or random GGP players were pitted against one another.

For the IGGP problem information about games is represented as sets of ground atoms. These atoms can be divided into two sets, the atoms that can change during the game and the atoms that can not. The set of atoms that can change can again be divided in two, the action atoms $A$ and the state atoms $S$.

The set $S$ represents all atoms that can change from one state to the other. Examples of elements of this set $S$ might be the state of the board in \textit{tic tac toe} along with which of the players is to take their turn next:
\begin{verbatim}
( control noughts )
( cell 1 1 b ) ( cell 1 2 x ) ( cell 1 3 o )
( cell 2 1 x ) ( cell 2 2 o ) ( cell 2 3 o )
( cell 3 1 b ) ( cell 3 2 x ) ( cell 3 3 b )
\end{verbatim}

The set $A$ is the set of ground atoms representing the actions that can be taken. The moves that each player makes are taken from this set. For a game such as \textit{eight puzzle} these would consist of the set of atoms:
\[\{\texttt{( move i j )}| 0<i\leq 9, 0<j\leq 9\}\]
Each atom represent sliding the tile at $(i,j)$ into the space on the board. To represent a player not making a move we also always have $noop \in A$.

\subsubsection{Game traces}

We modified the GGP matchmaker to log the following information:
\begin{itemize}
	\item \textbf{Game state trace} - The sequence of game states: $states = (s_1,...,s_n)$ where each $s_i \subseteq S$ is the set of ground atoms true in the $i^{th}$ state.
	\item \textbf{Game roles} - The list of roles of each player in the game: $roles = (r_1,...,r_k)$ e.g. \textit{noughts} and \textit{crosses} in \textit{tic tac toe}
	\item \textbf{Move trace} - The sequence of moves made by each player after each state: $moves = ((m_{1,r_1},...,m_{1,r_k}),...,(m_{n,r_1},...,m_{n,r_k}))$ where $m_{i,r_j} \in A$ is the $i^{th}$ move of player $r_j$. In games where not all players move every turn then $m_{i,r_j}=noop$ shows that $r_j$ made no move
	\item \textbf{Legal move trace} - The sequence of the legal moves for each player in each state:  $legal = ((l_{1,r_1},...,l_{1,r_k}),...,(l_{n,r_1},...,l_{n,r_k}))$. $l_{i,r_j} \subseteq A$ is the set of possible moves for $r_j$ in state $s_i$
	\item \textbf{Goal value trace} - The sequence of goal value for each player on every state: $goals = ((g_{1,r_1},...,g_{1,r_k}),...,(g_{n,r_1},...,g_{n,r_k}))$. $g_{i,r_j} \in [0,100]$ represents how well player $r_j$ has achieved the goal in state $s_i$.

\end{itemize}
This represents all the data needed from the match to generate IGGP tasks. In each state all the positive atoms in $S$ (that is the atoms true in the current state) are arguments of the \texttt{true} predicate. For example if at the start of the game the first cell is blank on the board then we would have $\texttt{( true ( cell 1 1 b ) )} \in s_1$. All the positive atoms in $A$ are arguments of the \texttt{does} predicate in a similar fashion.

\subsection{Transforming game traces into IGGP tasks}

We define a function that transforms the sequences into four separate induction tasks. An induction task is the set of triples $\{(B_i,E_i^+,E_i^-)\}_{i=0}^n$. We first define a function trace that translates the sequences given in the log of the match into a set of pairs $\{(B_i,E_i^+)\}_{i=0}^n$ . We then define a function \textit{triple} which gives the the set of triples  $\{(B_i,E_i^+,E_i^-)\}_{i=0}^n$ from the set of pairs $\{(B_i,E_i^+)\}_{i=0}^n$.
\\

We flatten the sequences of $(m_{i,r_1},...,m_{i,r_k})$ in \textit{moves}, $(l_{1,r_1},...,l_{1,r_k})$ in \textit{legals} and $(g_{1,r_1},...,g_{1,r_k})$ in \textit{goals} to a set that includes the role name as an extra argument of the predicate (the arity is increased by 1). For example if
\[(m_{i,r_1},m_{i,r_2}) = \texttt{[( move 2 2 ), ( move 2 3 )]}\ \ \ \ \ \ \ \ \ \ \ \ \ \ \ \ \ \ \ \ \ \ \]
and $r_1 = black$ and $r_2 = white$ then it would be replaced by $m_i$
\[m_i = \texttt{\{( move black 2 2 ), ( move white 2 3 )\}}\ \ \ \ \ \ \ \ \ \ \ \ \ \ \ \ \]
\textit{moves} is now a sequence $(m_1,...,m_n)$. With \textit{legals} and \textit{goals} a similar intuitive flattening procedure is applied. All four predicates are now represented by a single flat list.
\subsubsection{Building the traces}
The trace function is defined for \textit{legal}, \textit{goal}, \textit{next} and \textit{terminal}. We treat the output of the zip function, which takes a pair of lists and turns them into a list of pairs, as a set rather than a list. We define $S[p/q]$ on a set $S$ for predicates $p$ and $q$ to be the substitution of all instances of $q$ in $S$ for $p$.
\begin{align*}
&\Lambda_{legal} &&= trace_{legal}(states,legal) &&= \text{zip } states\ legal[\texttt{legal}/\texttt{true}] \\
&\Lambda_{goal} &&= trace_{goal}(states,goals) &&= \text{zip \textit{states goal}} \\
&\Lambda_{next} &&= trace_{next}(states,moves) &&= \text{zip \textit{states} } \text{(tail }states[\texttt{next}/\texttt{true}]) \\
&\Lambda_{terminal} &&= trace_{terminal}(states) &&= \text{zip \textit{states t}} \\
&&&\ \ \ \ \textbf{where  } t =\text{take } n\ (\text{repeat } \emptyset ) &&\text{++}\ \ [\{\texttt{( terminal )}\}] \\
&&&\ \ \ \ \textbf{and } n = (\text{length }states) - 1 &&
\end{align*}

The substitutions $[\texttt{legal}/\texttt{true}]$ and $[\texttt{next}/\texttt{true}]$ ensure that all positive examples in $E^+$ are in the relevant predicate for training. Before, the list \textit{legal} consisted of the only atoms in the predicate \texttt{true}. For the ILP systems to identify this as an example of legal we need to surround these with the \texttt{legal} predicate. That is, we need all the \texttt{move} atoms in $legal$ to be inside the \texttt{legal} predicate, e.g.:
\[\texttt{( legal ( move 1 1 ) )}\]
Since the predicates are already in the \texttt{true} predicate we only need to substitute one out for the other.

Some of the ILP systems being tested cannot work with function symbols of arity greater than 0. To allow them to operate we merge the function and their arguments into one predicate. For example \texttt{( legal ( move 1 1 ) )} becomes \verb|( legal_move 1 1 )| where \verb|legal_move| is a newly formed predicate.

To generate the triples we use the two functions $triples_1$ and $triples_2$. Let $pred\ X\ p$ be the subset of atoms in the set $X$ that use the predicate $p$.
\begin{align*}
tri&ple_1(\Lambda_p) = \text{map } f\ \Lambda\\
&\textbf{where } f\ (B,E^+) = (B,E^+,(pred\ S\ p) - E^+) \ \ \ \ \ \ \ \ \ \ \ \ \ \ \ \ \ \ \ \ \ \ \ \ \ \ \ \ \ \ \ \ \ \ \ \ \ \\
tri&ple_2(\Lambda_p) = \text{map } f\ \Lambda \\
&\textbf{where } f\ (B,E^+) = (B,E^+,(pred\ A\ p) - E^+)
\end{align*}
We generate the IGGP task with $triples_1$ and $triples_2$ as below:
\[\Delta = triples_1(\Lambda_{goal}) \cup triples_1(\Lambda_{terminal}) \cup triples_2(\Lambda_{next}) \cup triples_2(\Lambda_{legal})\]
The IGGP task $\Delta$ is set to the ILP systems as described in chapter \ref{ch:methodology}.
